\scanpage{179}

D'autre part $\beta$ est un mon. sous \unsure{la s\textsuperscript{l}|la seule} hypothèse $\underline H^0_{X/S}(\underline O_X)=\underline O_S$ ``universelle'' ; et il est surjectif donc bijectif si on suppose \uline{de plus} que \deleted{\unsure{$\mathrm{Pic}$}} $X$ a localmt une section sur $S$.

\uline{Exemple} où $\beta$ n'est pas surjectif : variété de \unsure{Sévéri|Sé véri|Severi}-Brauer non triviale sur un corps, \ldots\ On verra qu'il y a une \inserted{\unsure{théorie}} d'obstruction à la surjectivité de $\beta\alpha$, à valeurs dans un $H^2(S,G_m)$ \ldots.

\uline{Remarques}

On voit donc que le foncteur \uuline{Prépic} n'est pas représentable \deleted{en général, que \unsure{dans un}}\inserted{sauf dans des cas} \deleted{\unsure{pas|par}} triviaux sans intérêt ; $\mathrm{Prépic}'_{X/S}$ ne peut être représentable que si $\beta$ est un isom et si $\text{\uuline{Pic}}_{X/S}$ est représentable, \deleted{\unsure{(que} \unsure{si} $\unsure{g}$} [ce qui en pratique joint à des conditions comme celles déjà envisagée, y inclus l'existence d'une section de $X$ sur $S$].

Quant à $\text{\uuline{Pic}}_{X/S}$, il est représentable dans des cas \unsure{généraux|suffisamment} importants qui seront déterminés par la suite. C'est par suite lui le plus important pour notre propos.

% archive sheet number: 178
